\chapter{Curvature}

\begin{center}\itshape
An ant living on a surface cannot see the surface bend in space, but it can measure: lengths, angles, areas, the circumference of small circles. Gauss
discovered that this is enough to determine how curved the surface is. Integrate that curvature over a closed surface and the answer is $2\pi$ times a
whole number that does not depend on the shape at all, only on the number of holes. That is the Gauss--Bonnet theorem, the first and most beautiful
bridge from geometry to topology.
\end{center}

From 1818 to 1832, Carl Friedrich Gauss, professor of astronomy and the most famous mathematician in Europe, spent his summers measuring the Kingdom of
Hanover. He walked and rode across the country with theodolites, chose hilltops as stations, and invented an instrument, the heliotrope, that reflected
sunlight from one station to another over a hundred kilometres away. The survey was a chain of triangles, and its computation needed geometry on the
curved surface of the Earth.

Out of it came, in 1827, the \emph{Disquisitiones generales circa superficies curvas}, the founding text of differential geometry. Its central result is
a statement Gauss himself called the \emph{Theorema Egregium}, the ``remarkable theorem'': the curvature of a surface, defined by how its normal turns in
space, can be computed from measurements made within the surface alone. A sheet of paper can be rolled into a cylinder, because both have curvature zero;
no map of the Earth can be drawn without distortion, because the sphere has positive curvature and the plane has none.

A famous story says that Gauss measured the triangle formed by three mountain tops, the Brocken, the Hohehagen and the Inselsberg, with sides of about
$69$, $85$ and $107$ kilometres, to test whether the angles of a large triangle in physical space add up to two right angles. Gauss did measure that
triangle, and did think about the possibility that physical space is not Euclidean. But whether this measurement was meant as such a test is disputed:
A.~I.~Miller argued in 1972 that the story is a myth, and later historians have read the evidence differently. What is not disputed is that the excess
of that triangle over $180^\circ$, about fifteen seconds of arc, is exactly what the curvature of the Earth predicts, through a formula of this chapter.

We define curvature intrinsically, with Cartan's moving frames, prove the Theorema Egregium by showing that it agrees with the curvature of a surface in
space, and then prove the Gauss--Bonnet theorem for every compact oriented surface by Shiing-Shen Chern's argument of 1944: the Euler characteristic is
defined as the sum of the indices of a vector field, and the curvature integral proves that this sum does not depend on the field. We end with the
Laplacian of a surface, whose eigenvalues on a flat torus count sums of two squares.

\section{Riemannian metrics}

\begin{definition}
A \emph{Riemannian metric} on a manifold $M$ is a choice of inner product $\langle\cdot,\cdot\rangle_p$ on each tangent space $T_pM$, smooth in the sense that in
every chart the functions $g_{ij}=\langle\frac\partial{\partial x^i},\frac\partial{\partial x^j}\rangle$ are smooth. The \emph{length} of a piecewise smooth curve
$\gamma\colon[a,b]\to M$ is $\int_a^b|\gamma'(t)|\dd t$. A diffeomorphism $F$ between Riemannian manifolds is an \emph{isometry} if $\langle dF_pv,dF_pw\rangle=\langle v,w\rangle$ for all
$p,v,w$.
\end{definition}

A submanifold of $\R^N$ inherits a metric by restricting the Euclidean inner product to its tangent spaces; this is the metric the ant measures with. In a
chart it reads $\sum_{i,j}g_{ij}\dd x^i\dd x^j$, the classical ``line element''. For the unit sphere in spherical coordinates
$(\phi,\vartheta)\mapsto(\sin\phi\cos\vartheta,\sin\phi\sin\vartheta,\cos\phi)$, it is $\dd\phi^2+\sin^2\phi\,\dd\vartheta^2$. Every manifold has a Riemannian metric: in each chart
of a cover take the Euclidean metric of the chart, and glue with a partition of unity, $\langle v,w\rangle=\sum_j\rho_j\langle v,w\rangle_j$, a locally finite sum of positive
semidefinite forms which is positive definite because at each point some $\rho_j>0$.

\begin{example}[the hyperbolic half-plane]\label{ex:halfplane}
On $\mathbb H=\{(x,y):y>0\}$ the metric $\frac{\dd x^2+\dd y^2}{y^2}$ makes vertical lengths shrink as one rises: the vertical segment from $(0,1)$ to $(0,\varepsilon)$
has length $\int_\varepsilon^1\frac{\dd y}y=\log\frac1\varepsilon$, which tends to infinity, so the boundary line is infinitely far away. The maps
$z\mapsto\frac{az+b}{cz+d}$ with $a,b,c,d$ real and $ad-bc=1$ are isometries. Chapter~9 is devoted to this surface; we will compute below that its curvature
is $-1$ everywhere.
\end{example}

\section{Frames and the connection form}

From now on $M$ is an oriented surface with a Riemannian metric. An \emph{oriented orthonormal frame} on an open set $U$ is a pair of vector fields $(e_1,e_2)$,
orthonormal at each point and positively oriented. Applying the Gram--Schmidt process to the coordinate fields of an oriented chart produces one, so frames
exist locally. The \emph{dual coframe} $(\theta^1,\theta^2)$ is the pair of $1$-forms with $\theta^i(e_j)=\delta_{ij}$, and $\dd A=\theta^1\wedge\theta^2$ is the \emph{area form},
which does not depend on the frame by Proposition~\ref{prop:wedgedet} (a rotation has determinant $1$).

\begin{proposition}[Cartan's first structure equations]\label{prop:structure}
For an oriented orthonormal coframe $(\theta^1,\theta^2)$ there is a unique $1$-form $\omega$, the \emph{connection form}, with
\[
\dd\theta^1=\omega\wedge\theta^2,\qquad\dd\theta^2=-\omega\wedge\theta^1.
\]
If $(\tilde e_1,\tilde e_2)$ is another oriented orthonormal frame, $\tilde e_1=\cos\psi\,e_1+\sin\psi\,e_2$ and $\tilde e_2=-\sin\psi\,e_1+\cos\psi\,e_2$ for a smooth local angle
function $\psi$, then $\tilde\omega=\omega+\dd\psi$.
\end{proposition}
\begin{proof}
Write $\dd\theta^1=a\,\theta^1\wedge\theta^2$ and $\dd\theta^2=b\,\theta^1\wedge\theta^2$ and put $\omega=a\theta^1+b\theta^2$. Then $\omega\wedge\theta^2=a\,\theta^1\wedge\theta^2$ and $-\omega\wedge\theta^1=b\,\theta^1\wedge\theta^2$.
If $\omega'$ also works, $(\omega-\omega')\wedge\theta^2=0$ and $(\omega-\omega')\wedge\theta^1=0$; writing $\omega-\omega'=c_1\theta^1+c_2\theta^2$, these say $c_1=0$ and $c_2=0$.

For the second statement, the dual coframe rotates in the same way, $\tilde\theta^1=\cos\psi\,\theta^1+\sin\psi\,\theta^2$ and $\tilde\theta^2=-\sin\psi\,\theta^1+\cos\psi\,\theta^2$. Then
\[
\dd\tilde\theta^1=\dd\psi\wedge(-\sin\psi\,\theta^1+\cos\psi\,\theta^2)+\cos\psi\,\omega\wedge\theta^2-\sin\psi\,\omega\wedge\theta^1=(\omega+\dd\psi)\wedge\tilde\theta^2,
\]
and similarly $\dd\tilde\theta^2=-(\omega+\dd\psi)\wedge\tilde\theta^1$. By uniqueness, $\tilde\omega=\omega+\dd\psi$. (The angle $\psi$ exists locally, as a continuous argument of the
component vector of $\tilde e_1$ in the frame $e$, by the lifting lemma of Tome~7, Lemma~9.1; it is smooth because locally it is the composite of that vector with a
branch of the argument.)
\end{proof}

\begin{definition}
The \emph{Gauss curvature} $K$ is the function defined by $\dd\omega=-K\,\dd A$.
\end{definition}

Since $\dd\tilde\omega=\dd\omega+\dd\dd\psi=\dd\omega$ and $\dd A$ does not depend on the frame, $K$ is a well defined smooth function on $M$, and it depends only on the metric.
If the orientation is reversed, both $\omega$ and $\dd A$ change sign (replace $e_2$ by $-e_2$), so $K$ does not even depend on the orientation.

\begin{proposition}[orthogonal coordinates]\label{prop:orthocurv}
If in a chart the metric is $E\,\dd x^2+G\,\dd y^2$, with $a=\sqrt E$ and $b=\sqrt G$, then
\[
\omega=-\frac{a_y}b\dd x+\frac{b_x}a\dd y,\qquad K=-\frac1{ab}\Bigl[\Bigl(\frac{b_x}a\Bigr)_x+\Bigl(\frac{a_y}b\Bigr)_y\Bigr].
\]
\end{proposition}
\begin{proof}
The coframe is $\theta^1=a\,\dd x$, $\theta^2=b\,\dd y$, with $\dd\theta^1=-a_y\dd x\wedge\dd y$ and $\dd\theta^2=b_x\dd x\wedge\dd y$. With $\omega=p\,\dd x+q\,\dd y$, the structure equations read
$pb=-a_y$ and $qa=b_x$. Then $\dd\omega=\bigl[(\frac{b_x}a)_x+(\frac{a_y}b)_y\bigr]\dd x\wedge\dd y$, and $\dd A=ab\,\dd x\wedge\dd y$.
\end{proof}

\begin{example}\label{ex:curvexamples}
For the sphere, $\dd\phi^2+\sin^2\phi\,\dd\vartheta^2$ gives $a=1$, $b=\sin\phi$, and $K=-\frac1{\sin\phi}(\cos\phi)_\phi=1$. For the hyperbolic half-plane, $a=b=\frac1y$ gives
$\frac{a_y}b=-\frac1y$, whose $y$-derivative is $\frac1{y^2}$, and $K=-y^2\cdot\frac1{y^2}=-1$. For the plane in polar coordinates, $\dd r^2+r^2\dd\vartheta^2$, $K=-\frac1r(1)_r=0$, as it should be.
\end{example}

\subsection*{Covariant derivative and geodesic curvature}

The connection form tells how the frame turns. For a vector field $Y=u_1e_1+u_2e_2$ and a tangent vector $v$, define
\[
\nabla_vY=\bigl(v(u_1)-u_2\,\omega(v)\bigr)e_1+\bigl(v(u_2)+u_1\,\omega(v)\bigr)e_2,
\]
so that $\nabla e_1=\omega\,e_2$ and $\nabla e_2=-\omega\,e_1$. By the rule $\tilde\omega=\omega+\dd\psi$ this does not depend on the frame: computing $\nabla\tilde e_1$ with the formula in the frame $e$ gives
$\dd\psi\,\tilde e_2+\omega\,\tilde e_2=\tilde\omega\,\tilde e_2$. This is the \emph{Levi-Civita covariant derivative} of the surface. The same formula differentiates vector fields
along a curve $\gamma$: $\frac{D}{\dd t}Y(t)=\nabla_{\gamma'(t)}Y$, with $v(u_i)$ replaced by $u_i'(t)$. A short computation shows
$\frac{\dd}{\dd t}\langle Y,Z\rangle=\langle\frac{DY}{\dd t},Z\rangle+\langle Y,\frac{DZ}{\dd t}\rangle$: the derivative respects the metric.

\begin{definition}
Let $\gamma$ be a curve parametrized by arc length, with unit tangent $T=\gamma'$, and let $JT$ be $T$ turned by $+90^\circ$ (so $(T,JT)$ is an oriented
orthonormal basis). Since $|T|=1$, $\frac{DT}{\dd s}$ is orthogonal to $T$, and the \emph{geodesic curvature} is the function $\kappa_g$ with
$\frac{DT}{\dd s}=\kappa_g\,JT$. A \emph{geodesic} is a curve with $\frac{D\gamma'}{\dd t}=0$.
\end{definition}

\begin{proposition}\label{prop:kappag}
If $T=\cos\varphi\,e_1+\sin\varphi\,e_2$ with $\varphi$ a continuous angle along $\gamma$, then $\kappa_g=\varphi'+\omega(T)$.
\end{proposition}
\begin{proof}
By the definition of $\nabla$ with $u_1=\cos\varphi$, $u_2=\sin\varphi$: $\frac{DT}{\dd s}=(-\varphi'\sin\varphi-\sin\varphi\,\omega(T))e_1+(\varphi'\cos\varphi+\cos\varphi\,\omega(T))e_2=(\varphi'+\omega(T))\,JT$.
\end{proof}

In a chart, the geodesic equation $\frac{D\gamma'}{\dd t}=0$ is a system of two second-order differential equations with smooth coefficients, so by Theorems~\ref{thm:picard}
and~\ref{thm:flowsmooth} for every $p$ and $v\in T_pM$ there is a unique geodesic $\gamma_v$ with $\gamma_v(0)=p$ and $\gamma_v'(0)=v$, depending smoothly on $(p,v)$. Its speed is
constant, since $\frac{\dd}{\dd t}|\gamma'|^2=2\langle\frac{D\gamma'}{\dd t},\gamma'\rangle=0$, and $\gamma_{cv}(t)=\gamma_v(ct)$ by uniqueness. The \emph{exponential map}
$\exp_p(v)=\gamma_v(1)$ is therefore defined and smooth for small $v$, with $\exp_p(tv)=\gamma_v(t)$; its differential at $0$ is the identity, so by the inverse function
theorem it is a diffeomorphism from a small disc of $T_pM$ onto a neighbourhood of $p$. Geodesics are the straight lines of the surface: Exercise~\ref{ex:geodlength}
shows that short geodesic arcs are the shortest curves between their endpoints.

\section{Surfaces in space and the Theorema Egregium}

Let $S\subset\R^3$ be an oriented surface with the induced metric. Its orientation gives a unit normal $N(p)$, the vector such that $(e_1,e_2,N)$ is a positive basis of
$\R^3$ for any oriented orthonormal frame of $T_pS$. The \emph{Gauss map} $N\colon S\to S^2$ records the direction of the surface. Since $|N|=1$, the derivative $dN_p$
maps $T_pS$ into $N(p)^\perp=T_pS$: it is a linear endomorphism of the tangent plane, and $-dN_p$ is the \emph{shape operator}.

\begin{theorem}[Gauss, Theorema Egregium]\label{thm:egregium}
For an oriented surface $S\subset\R^3$, $K=\det(dN)$. In particular $\det(dN)$, which is defined by how the surface sits in space, depends only on the metric of
$S$, and is preserved by isometries.
\end{theorem}
\begin{proof}
Work on a frame domain, with $e_3=N$, and regard $x$ (the position), $e_1,e_2,e_3$ as $\R^3$-valued functions on $S$. Define $1$-forms $\omega_{ij}=\langle\dd e_i,e_j\rangle$;
differentiating $\langle e_i,e_j\rangle=\delta_{ij}$ gives $\omega_{ij}=-\omega_{ji}$, and $\dd e_i=\sum_j\omega_{ij}e_j$. Also $\dd x=\theta^1e_1+\theta^2e_2$, since both sides take the value $v$ on
$v\in T_pS$.

Apply $\dd$ to $\dd x$, componentwise in $\R^3$: for a $1$-form $\theta$ and a vector function $e$, $\dd(\theta e)=\dd\theta\,e-\theta\wedge\dd e$. So
\[
0=\dd\theta^1e_1-\theta^1\wedge(\omega_{12}e_2+\omega_{13}e_3)+\dd\theta^2e_2-\theta^2\wedge(\omega_{21}e_1+\omega_{23}e_3).
\]
The $e_1$-component gives $\dd\theta^1=\theta^2\wedge\omega_{21}=\omega_{12}\wedge\theta^2$, the $e_2$-component $\dd\theta^2=\theta^1\wedge\omega_{12}=-\omega_{12}\wedge\theta^1$. By uniqueness in
Proposition~\ref{prop:structure}, $\omega=\omega_{12}$. (So $\nabla$ is the tangential part of the ordinary derivative in $\R^3$.) Now apply $\dd$ to $\dd e_1=\omega_{12}e_2+\omega_{13}e_3$ and take the
$e_2$-component: $0=\dd\omega_{12}-\omega_{13}\wedge\omega_{32}$, that is,
\[
\dd\omega=\omega_{13}\wedge\omega_{32}=-\omega_{31}\wedge\omega_{32}.
\]
Finally $\dd N=\dd e_3=\omega_{31}e_1+\omega_{32}e_2$, so the matrix of $dN$ in the basis $(e_1,e_2)$ has rows $(\omega_{31}(e_1),\omega_{31}(e_2))$ and $(\omega_{32}(e_1),\omega_{32}(e_2))$, and
by Proposition~\ref{prop:wedgedet}, $\omega_{31}\wedge\omega_{32}=\det(dN)\,\theta^1\wedge\theta^2$. Comparing with $\dd\omega=-K\,\dd A$ gives $K=\det(dN)$.
\end{proof}

The $e_3$-component of $\dd(\dd x)=0$ gives $\theta^1\wedge\omega_{13}+\theta^2\wedge\omega_{23}=0$, which says that the matrix of $dN$ is symmetric: the shape operator is self-adjoint. Its
eigenvalues $\kappa_1,\kappa_2$ are the \emph{principal curvatures} (up to sign), the maximum and minimum of the curvature of the normal sections, and
$K=\kappa_1\kappa_2$. The sphere of radius $R$ has $N(x)=\frac xR$, $dN=\frac1RI$ and $K=\frac1{R^2}$. A cylinder of radius $R$ has $\kappa_1=\frac1R$, $\kappa_2=0$ and $K=0$, and
indeed it is isometric to a strip of the plane, which one sees by unrolling it. The individual $\kappa_i$ are not intrinsic; only their product is.

\begin{beauty}
Gauss's proof was a page of formulas in which the second fundamental form, the extrinsic data, cancels miraculously from a certain combination. With
frames the miracle has a reason. The curvature is $\dd\omega$, and $\omega$ is determined by the coframe alone through $\dd\theta^1=\omega\wedge\theta^2$; the extrinsic forms
$\omega_{13},\omega_{23}$ enter only through the identity $\dd\dd e_1=0$, which forces $\omega_{13}\wedge\omega_{32}$, the determinant of the Gauss map, to equal something intrinsic.
\end{beauty}

\section{The Gauss--Bonnet theorem}

\subsection*{Indices of vector fields}

Let $X$ be a vector field on a surface with an isolated zero at $p$. In an oriented chart around $p$, $X$ becomes a map from a punctured disc to $\R^2\setminus\{0\}$, and the
\emph{index} $\operatorname{ind}_p(X)$ is the winding number around $0$ of the closed curve $t\mapsto X(c(t))$, where $c$ is a small circle around $p$ traversed anticlockwise.
Figure~\ref{fig:index} shows indices $+1$, $-1$ and $+2$.

\begin{lemma}\label{lem:index}
The index does not depend on the radius of the circle, nor on the oriented chart. It can also be computed in any oriented orthonormal frame $(f_1,f_2)$
defined on a disc around $p$: if $X=|X|(\cos\psi f_1+\sin\psi f_2)$, then $\operatorname{ind}_p(X)=\frac1{2\pi}\int_c\dd\psi$.
\end{lemma}
\begin{proof}
All three statements follow from the homotopy invariance of the winding number (Tome~7, Proposition~9.3). Shrinking the circle is a homotopy of closed curves
avoiding $0$. A change of oriented chart multiplies the components of $X$ by the Jacobian matrix $A(q)$ of the transition map, which has positive determinant;
along the circle of radius $r$, $s\mapsto A(rs\cdot)$, $s\in[0,1]$, deforms $A$ to the constant matrix $A(p)$, and the winding number of $t\mapsto A(p)u(t)$ equals that
of $u$, since $A(p)$ can be joined to the identity inside $\mathrm{GL}_2^+(\R)$ (write $A(p)=QR$ with $Q$ a rotation and $R$ upper triangular with positive diagonal, and
deform $R$ linearly to $I$ and $Q$ through rotations). For the frame: the components of $X$ in the frame $f$ are obtained from its chart components by a matrix
with positive determinant, so the same argument applies, and the winding number of the curve of frame components is $\frac1{2\pi}\int_c\dd\psi$ (Corollary~\ref{cor:windingintegral}).
\end{proof}

\begin{figure}[ht]
\centering
\fig{ch04-index}
\caption{Three isolated zeros. Going once around the orange circle anticlockwise, the arrow turns once anticlockwise for the source $(x,y)$, once
clockwise for the saddle $(x,-y)$, and twice anticlockwise for the field $z^2=(x^2-y^2,2xy)$. Their indices are $+1$, $-1$ and $+2$.}
\label{fig:index}
\end{figure}

\begin{theorem}[Gauss--Bonnet--Chern]\label{thm:gaussbonnet}
Let $M$ be a compact oriented surface with a Riemannian metric and $X$ a vector field on $M$ with finitely many zeros $p_1,\dots,p_m$. Then
\[
\int_MK\,\dd A=2\pi\sum_{i=1}^m\operatorname{ind}_{p_i}(X).
\]
\end{theorem}
\begin{proof}
On $M^*=M\setminus\{p_1,\dots,p_m\}$ let $e_1=\frac X{|X|}$ and $e_2=Je_1$. This is an oriented orthonormal frame on all of $M^*$, so its connection form $\omega$ is a $1$-form
on $M^*$ with $\dd\omega=-K\dd A$. Around each $p_i$ take a chart and the closed disc $D_i(\varepsilon)$ of radius $\varepsilon$ in it, and let $M_\varepsilon=M\setminus\bigcup_i\operatorname{int}D_i(\varepsilon)$, a compact
surface with boundary. Its boundary is the union of the circles $\partial D_i$, with the orientation opposite to the one they get as boundaries of the discs.
By Stokes' theorem,
\[
\int_{M_\varepsilon}K\dd A=-\int_{M_\varepsilon}\dd\omega=\sum_i\int_{\partial D_i(\varepsilon)}\omega,
\]
the circles now oriented as boundaries of the discs. Near $p_i$ choose an oriented orthonormal frame $(f_1,f_2)$ on a whole disc around $p_i$, with connection form
$\omega_f$, smooth across $p_i$, and write $e_1=\cos\psi f_1+\sin\psi f_2$ with a local angle $\psi$. By Proposition~\ref{prop:structure}, $\omega=\omega_f+\dd\psi$, so
\[
\int_{\partial D_i(\varepsilon)}\omega=\int_{D_i(\varepsilon)}\dd\omega_f+\int_{\partial D_i(\varepsilon)}\dd\psi=-\int_{D_i(\varepsilon)}K\dd A+2\pi\operatorname{ind}_{p_i}(X)
\]
by Stokes on the disc and Lemma~\ref{lem:index}. As $\varepsilon\to0$, the integrals of $K$ over the shrinking discs tend to $0$, and $\int_{M_\varepsilon}K\dd A\to\int_MK\dd A$.
\end{proof}

The left side does not involve $X$ and the right side does not involve the metric. So the sum of the indices is the same for all vector fields with finitely many
zeros, and the integral of the curvature is the same for all metrics. (A diffeomorphism carries vector fields to vector fields with the same indices, so the
sum is also a diffeomorphism invariant.)

\begin{definition}
The \emph{Euler characteristic} $\chi(M)$ of a compact oriented surface is $\sum_p\operatorname{ind}_p(X)$ for any vector field $X$ with finitely many zeros.
\end{definition}

For this definition to apply, such vector fields must exist; we show next that they always do, using Sard's theorem. Granting that, we have proved:

\begin{corollary}[Gauss--Bonnet]\label{cor:gaussbonnet}
For every Riemannian metric on a compact oriented surface $M$, $\int_MK\dd A=2\pi\chi(M)$.
\end{corollary}

\subsection*{Morse functions exist}

A critical point $p$ of a smooth function $f$ is \emph{non-degenerate} if the Hessian matrix $\bigl(\partial_i\partial_jf(p)\bigr)$ in a chart is invertible; this does not depend on the
chart, since at a critical point the Hessian transforms as $H\mapsto A^\top HA$ with $A$ the Jacobian of the transition map. A function whose critical points are all
non-degenerate is a \emph{Morse function}.

\begin{lemma}\label{lem:morseindex}
Let $f$ be a smooth function on a Riemannian surface and $p$ a non-degenerate critical point. Then $p$ is an isolated zero of $\operatorname{grad}f$ (the vector field with
$\langle\operatorname{grad}f,v\rangle=\dd f(v)$), and its index is the sign of $\det$ of the Hessian: $+1$ at a local maximum or minimum, $-1$ at a saddle.
\end{lemma}
\begin{proof}
In a chart with $p=0$, $\operatorname{grad}f=G^{-1}\nabla f$ where $G=(g_{ij})$ is positive definite and $\nabla f$ is the vector of partial derivatives. Since $D(\nabla f)(0)=H$ is invertible,
$\nabla f$ is a local diffeomorphism at $0$ by the inverse function theorem, so it does not vanish on a punctured neighbourhood: the zero is isolated. On a small
circle, $\nabla f(x)=Hx+o(|x|)$, so by the dog-on-a-leash principle (Tome~7, Proposition~9.3(a)) the winding numbers of $G^{-1}\nabla f$, of $\nabla f$ (deform $G^{-1}$ linearly to $I$
through positive definite matrices) and of $x\mapsto Hx$ agree. The last is $+1$ if $\det H>0$ and $-1$ if $\det H<0$: join $H$ inside $\mathrm{GL}_2(\R)$ to $I$ or to
$\operatorname{diag}(1,-1)$ as in Lemma~\ref{lem:index}. At a maximum or minimum $H$ is definite, so $\det H>0$; at a saddle it is indefinite, so $\det H<0$.
\end{proof}

\begin{theorem}\label{thm:morseexists}
Every compact surface $M$ carries a Morse function. Consequently it carries a vector field with finitely many zeros, and $\chi(M)$ is defined.
\end{theorem}
\begin{proof}
By Whitney's theorem (Theorem~\ref{thm:whitney}) we may take $M\subset\R^N$. For $a\in\R^N$ let $f_a(x)=\langle a,x\rangle$ on $M$. We show that $f_a$ is a Morse function for almost every $a$.
Near each point, after permuting the coordinates of $\R^N$, $M$ is a graph $x''=g(x')$ over an open set $\Omega\subset\R^2$ (Theorem~\ref{thm:fourdesc}(c)); the map $x'\mapsto(x',g(x'))$ is a chart.
Write $a=(a',a'')$. In this chart $f_a(x')=\langle a',x'\rangle+\langle a'',g(x')\rangle$, with gradient $a'+\nabla F_{a''}(x')$, where $F_{a''}=\langle a'',g\rangle$, and Hessian
$\operatorname{Hess}F_{a''}(x')$. So $x'$ is a degenerate critical point of $f_a$ exactly when $-a'=\nabla F_{a''}(x')$ is a critical value of the map
$\nabla F_{a''}\colon\Omega\to\R^2$. For fixed $a''$ these values form a null set of $a'\in\R^2$, by Sard's theorem in equal dimensions (Theorem~\ref{thm:sard}). By Fubini's theorem
(Tome~6), the set of $a\in\R^N$ for which $f_a$ has a degenerate critical point in this chart is null. (It is measurable: it is the image under the continuous map
$(x',a'')\mapsto(-\nabla F_{a''}(x'),a'')$ of the closed set where the Hessian is singular, a countable union of compact sets.) Countably many such charts cover $M$, so
for almost every $a$, $f_a$ is a Morse function. Its critical points are isolated by Lemma~\ref{lem:morseindex}, hence finite in number since $M$ is compact, and they are
the zeros of $\operatorname{grad}f_a$ for any Riemannian metric on $M$.
\end{proof}

\begin{corollary}\label{cor:chimorse}
If $f$ is a Morse function on a compact oriented surface, $\chi(M)=\#\{\text{minima}\}-\#\{\text{saddles}\}+\#\{\text{maxima}\}$.
\end{corollary}

\begin{example}\label{ex:chiexamples}
The height function $z$ on $S^2$ has one minimum and one maximum, so $\chi(S^2)=2$, and Gauss--Bonnet with $K=1$ gives $4\pi$, the area of the sphere. The torus
carries the nowhere-vanishing field $\partial_{\theta_1}$, so $\chi(T^2)=0$: every metric on a torus has curvature of both signs, or is flat. For the surface of genus $g$
standing upright in space, like a pretzel with $g$ holes, the height function has one minimum, one maximum and two saddles for each hole (at the bottom and
the top of the hole), so $\chi=2-2g$. Chapter~6 shows that the alternating count $V-E+F$ is the same for every triangulation of a surface, and that it is also $2-2g$ for these surfaces.
\end{example}

\begin{remark}
Theorem~\ref{thm:gaussbonnet} with $M=S^2$ reproves the hairy ball theorem in dimension $2$: a vector field without zeros would give $\chi(S^2)=0$. The analogous
statement in every dimension, $\sum\operatorname{ind}(X)=\chi(M)$ for any vector field with isolated zeros on a compact manifold, is the Poincar\'e--Hopf theorem (Poincar\'e,
1885, for surfaces; Hopf, 1926); Chern's 1944 paper proves the corresponding Gauss--Bonnet formula in every even dimension by the same transgression on the unit
tangent bundle.
\end{remark}

\section{The Laplacian of a surface}

On an oriented Riemannian surface, the \emph{Hodge star} of $1$-forms is the rotation by $90^\circ$: $\ast\theta^1=\theta^2$, $\ast\theta^2=-\theta^1$ in any oriented orthonormal coframe (it
commutes with rotations, so it does not depend on the coframe). On $2$-forms, $\ast(h\,\dd A)=h$. The \emph{Laplace--Beltrami operator} is
\[
\Delta f=\ast\,\dd\ast\dd f.
\]
In the coframe $\dd x,\dd y$ of the Euclidean plane, $\ast\dd f=f_x\dd y-f_y\dd x$ and $\dd\ast\dd f=(f_{xx}+f_{yy})\dd x\wedge\dd y$: it is the ordinary Laplacian. The definition uses only
the metric, so $\Delta$ commutes with isometries.

\begin{proposition}[Green's formula]\label{prop:green}
On a compact oriented surface without boundary, $\int_Mg\,\Delta f\dd A=-\int_M\langle\dd g,\dd f\rangle\dd A$. Hence $-\Delta$ is symmetric and non-negative, its eigenvalues
($\Delta f=-\lambda f$, $f\ne0$) satisfy $\lambda\ge0$, and if $M$ is connected the eigenvalue $0$ has only the constants as eigenfunctions.
\end{proposition}
\begin{proof}
With $\dd f=f_1\theta^1+f_2\theta^2$ and $\dd g=g_1\theta^1+g_2\theta^2$, $\dd g\wedge\ast\dd f=(g_1f_1+g_2f_2)\dd A$. So $\dd(g\ast\dd f)=\langle\dd g,\dd f\rangle\dd A+g\,\Delta f\,\dd A$, and Stokes gives $0$ for the
integral of the left side. Taking $g=f$ for an eigenfunction gives $\lambda\int f^2=\int|\dd f|^2\ge0$, with equality only if $\dd f=0$.
\end{proof}

\begin{theorem}[the spectrum of the square flat torus]\label{thm:torusspectrum}
On $T^2=\R^2/\Z^2$ with the flat metric $\dd x^2+\dd y^2$, the eigenvalues of $-\Delta$ are the numbers $4\pi^2n$ with $n=k_1^2+k_2^2$ a sum of two squares, and the multiplicity of
$4\pi^2n$ is $r_2(n)$, the number of $(k_1,k_2)\in\Z^2$ with $k_1^2+k_2^2=n$. The functions $e_k(x)=e^{2\pi i\langle k,x\rangle}$, $k\in\Z^2$, form an orthonormal basis of eigenfunctions of
$L^2(T^2)$.
\end{theorem}
\begin{proof}
$\Delta e_k=-4\pi^2|k|^2e_k$ by direct computation, and the $e_k$ are orthonormal by Fubini and the one-variable computation. They span a dense subspace: the
trigonometric polynomials form an algebra of continuous functions on the compact space $T^2$, stable under conjugation, containing the constants and separating
points, hence dense in $C(T^2)$ for the uniform norm by Stone--Weierstrass (Tome~6, Theorem~4.19, applied to real and imaginary parts), and $C(T^2)$ is dense in
$L^2$. Now let $\Delta f=-\lambda f$ with $f$ smooth. By Green's formula applied to real and imaginary parts, $\int\overline{e_k}\Delta f=\int\Delta\overline{e_k}\,f$, so
$-\lambda\hat f(k)=-4\pi^2|k|^2\hat f(k)$ for every Fourier coefficient. Hence $\hat f(k)=0$ unless $4\pi^2|k|^2=\lambda$, and $f$ is a combination of the $e_k$ with $|k|^2=\frac\lambda{4\pi^2}$.
\end{proof}

\begin{remark}
The counting function $N(\lambda)=\#\{\text{eigenvalues}\le\lambda\}$, with multiplicity, is the number of lattice points in the disc of radius $\frac{\sqrt\lambda}{2\pi}$, so
$N(\lambda)\sim\frac\lambda{4\pi}$: the area of the torus times $\frac\lambda{4\pi}$. Weyl proved in 1911 that this holds for every bounded plane domain, so that one ``hears'' the
area of a drum; how much more one can hear is Mark Kac's 1966 question ``Can one hear the shape of a drum?''. The error term is the Gauss circle problem of
Tome~5. On the round sphere the eigenvalues are $\ell(\ell+1)$ with multiplicity $2\ell+1$, the eigenfunctions being the spherical harmonics; we do not prove this here.
Milnor found in 1964 two flat tori of dimension $16$ that are not isometric but have the same spectrum: one cannot always hear the shape of a torus.
\end{remark}

\section{Exercises}

\begin{figure}[ht]
\centering
\fig{ch04-girard}
\caption{A spherical triangle $ABC$ on the unit sphere, with the three great circles that carry its sides. Its angles add up to more than $\pi$; the excess is its area.}
\label{fig:girard}
\end{figure}

\begin{exercise}[$\star$]\label{ex:girard}
\emph{Girard's theorem.} A \emph{lune} of angle $\alpha$ on the unit sphere is the region between two half great circles meeting at angle $\alpha$ at antipodal points. (a) Show that
its area is $2\alpha$. (b) Let $ABC$ be a spherical triangle with angles $\alpha,\beta,\gamma$ (Figure~\ref{fig:girard}). Extend its sides to full great circles and show that the three
lunes containing the triangle, together with their antipodal images, cover the sphere, covering the triangle and its antipode three times and every other point
once. (c) Deduce that the area of the triangle is $\alpha+\beta+\gamma-\pi$. (d) For the Brocken--Hohehagen--Inselsberg triangle, of area about $2.9\cdot10^3$ km$^2$ on an Earth of radius
$6371$ km, estimate the excess in seconds of arc.
\end{exercise}

\begin{exercise}[$\star$]\label{ex:revolution}
\emph{Surfaces of revolution.} Let $s\mapsto(r(s),z(s))$ be a curve in the half-plane $r>0$ parametrized by arc length, and $S$ the surface
$(r(s)\cos\vartheta,r(s)\sin\vartheta,z(s))$. (a) Show that the metric is $\dd s^2+r^2\dd\vartheta^2$ and $K=-\frac{r''}r$. (b) For the torus obtained by rotating the circle of radius $a$
centred at distance $R>a$ from the axis, compute $K$ and check $\int K\dd A=0$. Where is $K$ positive, and where negative?
\end{exercise}

\begin{exercise}[$\star\star$]\label{ex:pseudosphere}
\emph{The pseudosphere.} (a) Find all profiles with $K=-1$ and $r(s)=e^{-s}$, $s>0$, and show that the corresponding surface of revolution (the \emph{pseudosphere}, generated by a
tractrix) exists. (b) Show that $(s,\vartheta)\mapsto(\vartheta,e^s)$ is a local isometry from the pseudosphere onto the region $y>1$ of the hyperbolic half-plane, which identifies $x$ with $x+2\pi$.
(c) What is the area of the pseudosphere? Compare with a hemisphere of radius $1$.
\end{exercise}

\begin{exercise}[$\star$]\label{ex:chisphere}
\emph{Counting zeros on the sphere.} (a) Show that the fields $z^n$ and $\bar z^n$ on $\C$ (viewed as vector fields) have index $n$ and $-n$ at $0$. (b) Let $X$ be the constant field $\partial_x$ in the
stereographic chart $z$ of $S^2\setminus\{N\}$. Show that in the chart $w=\frac1z$ around $N$ it becomes $-w^2\partial_w$, and deduce again that $\chi(S^2)=2$.
\end{exercise}

\begin{exercise}[$\star$]\label{ex:positivepoint}
\emph{Compact surfaces in space are somewhere convex.} Let $S\subset\R^3$ be a compact surface and $p$ a point of $S$ farthest from the origin, at distance $R$. Show that
$K(p)\ge\frac1{R^2}$. Deduce that no compact surface in $\R^3$ has $K\le0$ everywhere, and that a torus in $\R^3$ has points of negative curvature.
\end{exercise}

\begin{exercise}[$\star\star$]\label{ex:geodlength}
\emph{Geodesics are locally shortest.} Let $p\in M$ and $U=\exp_p(D)$ with $D$ a small disc of radius $\rho$ on which $\exp_p$ is a diffeomorphism. (a) Using polar coordinates
$(r,\vartheta)$ on $D$ transported by $\exp_p$, admit the \emph{Gauss lemma}: the metric takes the form $\dd r^2+G(r,\vartheta)\dd\vartheta^2$. Show that every piecewise smooth curve from $p$ to
$\exp_p(v)$, $|v|<\rho$, has length at least $|v|$, with equality only for the geodesic. (b) Prove the Gauss lemma in the case of the sphere and of the half-plane by
computing $\exp_p$ explicitly at one point.
\end{exercise}

\section{Hints}

\begin{hint}{ex:girard}
The lunes of angle $\alpha$ through $A$ and its antipode $A'$ form a pair. Total area $4\pi$.
\end{hint}
\begin{hint}{ex:revolution}
Proposition~\ref{prop:orthocurv} with $a=1$, $b=r$. For the torus, $r=R+a\cos\frac sa$ and $\dd A=r\dd s\dd\vartheta$.
\end{hint}
\begin{hint}{ex:pseudosphere}
$r''=r$; the profile needs $r'^2+z'^2=1$, which forces $|r'|\le1$. For (b) compare $\frac{\dd x^2+\dd y^2}{y^2}$ with $x=\vartheta$, $y=e^s$.
\end{hint}
\begin{hint}{ex:chisphere}
For (b), $\frac{\dd w}{\dd t}=-\frac1{z^2}\frac{\dd z}{\dd t}$.
\end{hint}
\begin{hint}{ex:positivepoint}
Every normal section of $S$ at $p$ lies inside the sphere of radius $R$ and touches it at $p$, so it curves at least as much.
\end{hint}
\begin{hint}{ex:geodlength}
In the coordinates $(r,\vartheta)$, the length of a curve is at least $\int|r'|\dd t$.
\end{hint}

\section{Solutions}

\begin{solution}{ex:girard}
\idea Areas of lunes are proportional to their angles, and the three great circles cut the sphere into eight triangles which the lunes cover in a countable way.
\writeup (a) Rotation about the axis $AA'$ preserves area, and $n$ lunes of angle $\frac{2\pi}n$ tile the sphere, so the area of a lune is proportional to its angle for angles
$\frac{2\pi m}n$, hence for all angles by monotonicity: area $=\frac\alpha{2\pi}4\pi=2\alpha$. (b) The three great circles cut the sphere into eight triangles: $T=ABC$, its antipode
$T'$, and six others. The lune at $A$ with angle $\alpha$ is $T\cup T_a$, where $T_a$ is the triangle across the side $BC$; its antipodal lune is $T'\cup T_a'$. Similarly for $B$ and $C$.
The six lunes consist of $T$ and $T'$ three times each, and each of the six other triangles once, and these six are all the other triangles. (c) Summing areas:
$2(2\alpha+2\beta+2\gamma)=4\pi+2\cdot2\operatorname{area}(T)$, that is, $\operatorname{area}(T)=\alpha+\beta+\gamma-\pi$. This is Gauss--Bonnet for a geodesic triangle with $K=1$: the integral of the curvature
is the excess of the angle sum. (d) On a sphere of radius $R$ areas scale by $R^2$, so the excess is $\frac{\text{area}}{R^2}=\frac{2.9\cdot10^3}{6371^2}\approx7.2\cdot10^{-5}$ radians, about
$15$ seconds of arc, the fifteen seconds Gauss had to account for. (The area $2.9\cdot10^3$ km$^2$ follows from the side lengths by Heron's formula.)
\end{solution}

\begin{solution}{ex:revolution}
\idea Parallels are orthogonal to meridians, so the metric is diagonal and Proposition~\ref{prop:orthocurv} applies; the curvature is the acceleration of the radius.
\writeup (a) The coordinate vectors are
\[(r'\cos\vartheta,\,r'\sin\vartheta,\,z')\quad\text{and}\quad(-r\sin\vartheta,\,r\cos\vartheta,\,0),\]
of lengths $1$ and $r$, and they are orthogonal. With $a=1$, $b=r(s)$:
$K=-\frac1r(r')'=-\frac{r''}r$. (b) The circle has unit-speed parametrization $r=R+a\cos\frac sa$, $z=a\sin\frac sa$, so $K=\frac{\cos(s/a)}{a(R+a\cos(s/a))}$. It is positive on the outer half
($|s|<\frac{\pi a}2$), negative on the inner half, zero on the top and bottom circles. With $\dd A=r\dd s\dd\vartheta$, $\int K\dd A=2\pi\int_0^{2\pi a}\frac{\cos(s/a)}a\dd s=0$, as Gauss--Bonnet requires
for $\chi(T^2)=0$: the positive curvature outside is exactly compensated by the negative curvature inside.
\end{solution}

\begin{solution}{ex:pseudosphere}
\idea Constant curvature $-1$ means $r''=r$; choosing the decaying solution gives a horn whose metric is the hyperbolic one in disguise.
\writeup (a) $K=-1$ means $r''=r$, so $r=\alpha e^s+\beta e^{-s}$; the choice $r=e^{-s}$ gives $r'^2=e^{-2s}<1$ for $s>0$, so $z(s)=\int_0^s\sqrt{1-e^{-2u}}\dd u$ is defined and the
profile is a unit-speed curve. It is the tractrix: the segment of the tangent from the curve to the axis has constant length $1$. (b) The metric of the pseudosphere is
$\dd s^2+e^{-2s}\dd\vartheta^2$. With $x=\vartheta$, $y=e^s$, $\frac{\dd x^2+\dd y^2}{y^2}=e^{-2s}\dd\vartheta^2+\dd s^2$. The map is a local isometry from $\{s>0\}$ onto $\{y>1\}$, and $\vartheta$ and $\vartheta+2\pi$ give the same
point. So the pseudosphere is a piece of the hyperbolic plane rolled up, as a cylinder is a piece of the Euclidean plane rolled up. Hilbert proved in 1901 that the
whole hyperbolic plane cannot be realized as a smooth surface in $\R^3$ with the induced metric; the pseudosphere has an edge at $s=0$. (c)
$\int_0^{2\pi}\int_0^\infty e^{-s}\dd s\dd\vartheta=2\pi$, the area of a unit hemisphere. (Its volume is also half that of the unit ball, as Huygens observed.)
\end{solution}

\begin{solution}{ex:chisphere}
\idea The index of $z^n$ counts how many times $z^n$ turns while $z$ turns once; on the sphere, a constant field must have a zero at infinity of index $2$.
\writeup (a) On the circle $z=re^{it}$, $z^n=r^ne^{int}$ has argument $nt$, so it winds $n$ times; $\bar z^n=r^ne^{-int}$ winds $-n$ times. (b) The flow of $X$ is $z\mapsto z+t$. In the
chart $w=\frac1z$ at the north pole (a holomorphic, hence orientation-preserving, change of chart), the flow is $w\mapsto\frac1{1/w+t}=\frac w{1+tw}$, whose velocity at $t=0$ is $-w^2$. So $X$ extends smoothly to $S^2$ with a single
zero at $N$, where it is $-w^2$, of index $2$ (the constant $-1$ does not change the winding). Thus $\chi(S^2)=2$.
\end{solution}

\begin{solution}{ex:positivepoint}
\idea At a farthest point the surface lies inside a sphere that touches it, so it is curved at least as much as that sphere in every direction.
\writeup The function $h(x)=|x|^2$ on $S$ has a maximum at $p$, so $\dd h_p=0$, which says that $p\perp T_pS$: the normal at $p$ is $\pm\frac pR$. Take $N(p)=\frac pR$. For a unit-speed curve
$c$ in $S$ with $c(0)=p$, $c'(0)=v$, the function $|c(t)|^2$ has a maximum at $0$, so its second derivative $2|v|^2+2\langle c''(0),p\rangle$ is $\le0$, that is,
$\langle c''(0),N(p)\rangle\le-\frac1R$. Differentiating $\langle c',N\circ c\rangle=0$ gives $\langle c'',N\rangle=-\langle v,dN(v)\rangle$, so $\langle dN_p(v),v\rangle\ge\frac1R$ for every unit $v$: both eigenvalues of the
symmetric map $dN_p$ are at least $\frac1R$, and $K(p)=\det dN_p\ge\frac1{R^2}$. So $K\le0$ everywhere is impossible. A torus has $\int K\dd A=0$ for any metric, and $K>0$ at $p$, so $K<0$
somewhere.
\end{solution}

\begin{solution}{ex:geodlength}
\idea In geodesic polar coordinates the radial direction has unit length and is orthogonal to the circles, so any curve advances in $r$ at most as fast as its speed.
\writeup (a) Let $c\colon[0,1]\to M$ be piecewise smooth from $p$ to $q=\exp_p(v)$. If $c$ leaves $U$, it must cross every circle $r=\rho'<\rho$, so its length exceeds $\rho'$ for every
$\rho'<\rho$, hence is at least $\rho>|v|$. Otherwise, let $t_0$ be the last time with $c(t_0)=p$; on $]t_0,1]$, $c(t)=\exp_p(r(t)\xi(t))$ with $r>0$, and in the metric
$\dd r^2+G\dd\vartheta^2$, $|c'|\ge|r'|$. So the length is at least $\int_{t_0}^1r'\dd t=|v|$. Equality forces $\vartheta'=0$ (where $G>0$) and $r'\ge0$, so $c$ is a reparametrization of the radial
geodesic. (b) On the unit sphere at the north pole, $\exp_p$ maps the vector of length $r$ in direction $\vartheta$ to the point at colatitude $r$ and longitude $\vartheta$ (great circles
are the geodesics, as $\nabla$ is the tangential part of the derivative), and the metric is $\dd r^2+\sin^2r\dd\vartheta^2$: the Gauss lemma with $G=\sin^2r$. On the half-plane at $i$,
geodesics through $i$ are the vertical line and the half-circles centred on the real axis (Chapter~9), and in geodesic polar coordinates the metric turns out to be
$\dd r^2+\sinh^2r\,\dd\vartheta^2$; we leave this case to Chapter~9, where the isometry group makes the computation immediate.
\end{solution}

\section*{Sources and further reading}
\addcontentsline{toc}{section}{Sources and further reading}

C.~F.~Gauss, ``Disquisitiones generales circa superficies curvas'', \emph{Comment.\ Soc.\ Reg.\ Sci.\ Gotting.\ Recent.} 6 (1828), presented on 8 October 1827;
English translation by A.~Hiltebeitel and J.~Morehead, \emph{General Investigations of Curved Surfaces} (Princeton, 1902; Dover reprint). On the survey and the triangle,
A.~I.~Miller, ``The myth of Gauss' experiment on the Euclidean nature of physical space'', \emph{Isis} 63 (1972), 345--348, and E.~Breitenberger, ``Gauss's
geodesy and the axiom of parallels'', \emph{Arch.\ Hist.\ Exact Sci.} 31 (1984), 273--289, which reads the evidence differently. W.~K.~B\"uhler,
\emph{Gauss: A Biographical Study} (Springer, 1981), Chapter~8, describes the Hanover survey and the heliotrope.

O.~Bonnet, ``M\'emoire sur la th\'eorie g\'en\'erale des surfaces'', \emph{J.\ \'Ec.\ Polytech.} 19 (1848), 1--146. S.-S.~Chern, ``A simple intrinsic proof of
the Gauss--Bonnet formula for closed Riemannian manifolds'', \emph{Ann.\ of Math.} 45 (1944), 747--752. H.~Hopf, ``Vektorfelder in $n$-dimensionalen
Mannigfaltigkeiten'', \emph{Math.\ Ann.} 96 (1926), 225--250. H.~Weyl, ``Das asymptotische Verteilungsgesetz der Eigenwerte linearer partieller
Differentialgleichungen'', \emph{Math.\ Ann.} 71 (1912), 441--479 (announced in 1911); M.~Kac, ``Can one hear the shape of a drum?'', \emph{Amer.\ Math.\ Monthly}
73 (1966), 1--23; J.~Milnor, ``Eigenvalues of the Laplace operator on certain manifolds'', \emph{Proc.\ Nat.\ Acad.\ Sci.\ USA} 51 (1964), 542.

The treatment by moving frames follows B.~O'Neill, \emph{Elementary Differential Geometry}, 2nd ed. (Academic Press, 2006), Chapters~6--7, and M.~do~Carmo,
\emph{Differential Forms and Applications} (Springer, 1994); for the classical approach, do~Carmo's \emph{Differential Geometry of Curves and Surfaces}
(Prentice-Hall, 1976). The existence of Morse functions follows V.~Guillemin and A.~Pollack, \emph{Differential Topology} (Prentice-Hall, 1974), Chapter~1, \S7.
